\documentclass[10pt,a4paper]{amsart}
\usepackage[margin=19mm]{geometry}
\usepackage{amsmath,amssymb,mathtools}
\usepackage[scr=rsfs,cal=euler]{mathalfa}
\usepackage{xfrac}
\usepackage[unicode,hidelinks,bookmarks=false]{hyperref}
\newcommand{\ie}{i.e.\ }
\newcommand{\cf}{cf.\ }
\newcommand{\resp}{respectively\ }
\newcommand{\Subsection}{\subsection}
\providecommand{\nobreakdash}{\nobreak\hbox{-}}
\setlength{\emergencystretch}{2em}
\multlinegap0pt
\usepackage{bm}
\usepackage{bbm}
\usepackage{dsfont}
\usepackage{xspace}
\usepackage{cancel}
\usepackage{mathtools}
\usepackage{cleveref}
\usepackage{amsthm}
\makeatletter
\@ifundefined{theorem}{\newtheorem{theorem}{Theorem}}{}
\@ifundefined{construction}{\newtheorem{construction}[theorem]{Construction}}{}
\@ifundefined{lemma}{\newtheorem{lemma}[theorem]{Lemma}}{}
\makeatother
\DeclarePairedDelimiter{\abs}{\lvert}{\rvert}
\newcommand{\bx}{\mathbf{x}}
\newcommand{\by}{\mathbf{y}}
\renewcommand{\emptyset}{\varnothing}
\renewcommand{\geq}{\geqslant}
\renewcommand{\leq}{\leqslant}
\title[Asymmetric results about graph homomorphisms]{Asymmetric results about graph homomorphisms}
\author{Lior Gishboliner, Eoin Hurley,, Yuval Wigderson}
\date{Source: arXiv:2502.20278v3 (CC BY 4.0)}
\begin{document}
\maketitle

\noindent\emph{Source and licence.} Asymmetric results about graph homomorphisms. Version arXiv:2502.20278v3 (CC BY 4.0), distributed under the Creative Commons Attribution 4.0 International licence, as recorded in the arXiv licence metadata for this exact version and shown on the abstract page https://arxiv.org/abs/2502.20278v3. Full rights notice: licenses/arxiv-2502.20278-CC-BY-4.0.txt.\par\smallskip

\noindent\textit{Editorial scope.} Counted unit: the Lemma whose source label is \texttt{lem:high-girth hypergraph} in the article (source0.tex lines 838--840, source label \texttt{lem:high-girth hypergraph}), delivered together with the article's own complete original proof of it (source0.tex lines 843--879, one \texttt{proof} environment of 37 lines). No mathematics is written, repaired or re-derived: statement and proof are the article's own text line for line. The proof is detached from the statement in the source: the article states the result at lines 838--840 and prints its proof at lines 843--879, and the proof environment's own header names the statement, so the association is the article's own. Required context retained uncounted: con:HG (lines 313--324); prop:no-approx-hom (lines 366--371); thm:approx hom hardness (lines 385--389); lem:hyperforest (lines 833--835). Every definition, display and label of those lines is retained unchanged. Omitted from this package and not redistributed: 1--312, 325--365, 372--384, 390--832, 836--837, 841--842, 880--966. Nothing inside a retained excerpt is omitted or altered: every retained body is delivered exactly as the article's lines are written, so no omission receipt of any kind accompanies this package. Citations: the retained excerpts carry no citation command at all, so no bibliography entry of the article is transcribed and no reference is redistributed. Reference adaptation: none was needed and none was made; every reference inside the delivered unit resolves inside it, so no unresolved reference or citation is produced. Printed slips kept literal: no printed word, symbol, hypothesis or number of the counted statement or of its proof was altered. Layout and class: the article's own class is \texttt{amsart} with packages including xcolor, amsfonts, amsmath, amssymb, amsthm, bbm, mathtools, comment, enumitem, url, hyperref, cleveref, euscript, adjustbox; the wrapper is the lane's pinned \texttt{amsart} editorial wrapper at 10pt on A4, which restates the theorem-like environments the retained text needs and adds the article's own macro definitions that the retained text needs (\texttt{\textbackslash abs}, \texttt{\textbackslash bx}, \texttt{\textbackslash by}, \texttt{\textbackslash emptyset}, \texttt{\textbackslash geq}, \texttt{\textbackslash leq}), with extra wrapper packages (bm, bbm, dsfont, xspace, cancel, mathtools, cleveref). The delivered numbering is the unit's own. Assets: no retained span contains a figure environment, a graphics-inclusion command, a PDF-page inclusion, a table or a TikZ picture; no image file is copied into this package, the package declares no asset dependency and no image file is redistributed with it. Licence: version arXiv:2502.20278v3 of this article is distributed under the Creative Commons Attribution 4.0 International licence, as recorded in the arXiv licence metadata for this exact version and shown on the abstract page https://arxiv.org/abs/2502.20278v3. A full rights notice is delivered at licenses/arxiv-2502.20278-CC-BY-4.0.txt. Characters: every retained excerpt is pure ASCII, so the delivered engine, XeLaTeX with its default Latin Modern text font, reports no missing character.
\begin{construction}
\label{con:HG}
Let $F$ be a graph and let $\Delta$ be its maximum degree. 
For each vertex $v\in V(F)$, fix an ordering of the edges incident to $v$ and label them $\{(v,1),\dots,(v,\deg(v))\}$.\footnote{Thus each edge receives two labels (one from each endpoint), and knowing one of them is enough to identify the edge.} 
When considering a copy of $F$ in a graph $G$, we will (implicitly) fix an isomorphism from $F$ to the copy, so that the copy is endowed with the above edge-labelling. 

Let $G$ be an $n$-vertex graph where every edge is contained in a unique copy of $F$.
Let $F_1,\dots,F_m$ be the copies of $F$ in $G$.
The graph $G^{\star}$ is constructed as follows. The vertex set\footnote{We denote a vertex of $G^{\star}$ as $(v,\bx)$, where $v \in V(G)$ and $\bx \in \{1,\dots,\Delta\}^m$.} of $G^{\star}$ is $V(G) \times \{1,\dots,\Delta\}^m$.
For every edge $e=uv \in E(G)$, let $k \in [m]$ be the unique index with $e \in E(F_k)$, and suppose that $e$ has labels $(u,i)$ and $(v,j)$ (as an edge of $F_k$). Join all vertices $(u,\bx)$ and $(v,\by)$ such that $\bx_k = i$ and $\by_k = j$.   
These are the only edges in $G^{\star}$. 
\end{construction}

\begin{lemma}\label{prop:no-approx-hom}
For every graph $F$, there exists $c>0$ such that the following holds.
Let $G$ be an $n$-vertex graph where every edge is contained in a unique copy of $F$, and let $m$ denote the number of copies of $F$ in $G$.
Let $G^{\star}$ be given by \cref{con:HG}. 
Then for $\varepsilon < \frac{cm}{n^2}$, there is no $\varepsilon$-approximate homomorphism from $G^{\star}$ to an $F$-hom-free graph on at most $2^{c m/n}$ vertices.
\end{lemma}

\begin{theorem}\label{thm:approx hom hardness}
	Let $F,H$ be graphs, where $F$ is $2$-connected, and suppose that $H$ is not homomorphic to $T^{\star}$ for any $F$-forest $T$ with $T\rightarrow F$. 
	Then for every small enough $\varepsilon > 0$,
	$M_{F,H}(\varepsilon) \geq 2^{(1/\varepsilon)^c}$, where $c > 0$ depends only on $F,H$. 
\end{theorem}

	\begin{lemma}\label{lem:hyperforest}
		A hypergraph $\mathcal{G}$ is a hyperforest if and only if it has no Berge cycles. 
	\end{lemma}

	\begin{lemma}\label{lem:high-girth hypergraph}
		For every $f,g \geq 2$ and every $n\geq f$, there is an $n$-vertex $f$-uniform, $f$-partite hypergraph with $\Omega(n^{1+1/g})$ edges and no Berge cycle of length at most $g$.
	\end{lemma}

	\begin{proof}[Proof of \cref{thm:approx hom hardness}]
	Set $f \coloneqq \abs F$, $h \coloneqq \abs H$.
	Let $\mathcal{G}$ be the hypergraph given by \cref{lem:high-girth hypergraph} with parameters $g \coloneqq \max(f,h)$ and $n \coloneqq c_0\cdot (\frac{1}{\varepsilon})^{\frac{1}{1-1/g}}$, where $c_0 > 0$ is a small enough constant to be chosen later. 
	{Let $G$ be the graph obtained from $\mathcal{G}$ by identifying the parts of $\mathcal{G}$ with the vertices of $F$, and accordingly placing a copy of $F$ on each hyperedge of $\mathcal{G}$, so that $G$ is a subgraph of a blowup of $F$.} 
	Letting $m$ denote the number of $F$-copies in $G$, we have $m \geq e(\mathcal G)= \Omega(n^{1 + 1/g})$. 
	Let $c = c(F)$ be the constant given by \cref{prop:no-approx-hom}.
	Note that 
	$$
	\frac{cm}{n^2} = \frac{\Omega(n^{1+1/g})}{n^2} = \Omega(n^{-1+1/g}) = \Omega( c_0^{-1+1/g} \cdot \varepsilon ) > \varepsilon,
	$$
	provided that $c_0$ is small enough. 
    Also, by definition, every edge of $G$ is on a copy of $F$. 


    Consider any set $X \subseteq V(G) = V(\mathcal{G})$ with $|X| \leq g$, and let $e_1,\dots,e_s$ be the hyperedges $e \in E(\mathcal{G})$ satisfying $|e \cap X| \geq 2$.  Consider the hypergraph 
    $\mathcal{G}_X = \{e_i \cap X : 1 \leq i \leq s\}$. Since $|X| \leq g$ and $\mathcal{G}$ has no Berge cycle of length at most $g$, we get that $\mathcal{G}_X$ has no Berge cycles. Therefore, $\mathcal{G}_X$ is a hyperforest by \cref{lem:hyperforest}. 
    In the case $|X| = f$, the graph $G[X]$ can contain a copy of $F$ only if $\mathcal{G}_X$ has a hyperedge of size $f$, because $F$ is 2-connected.
    Thus, the only $F$-copies in $G$ are those corresponding to hyperedges of $\mathcal{G}$. 
    As $\mathcal{G}$ has no Berge 2-cycles, it follows that every edge of $G$ is on a unique copy of $F$. This in turn allows us to apply \cref{con:HG}. 
    
	So let $G^{\star}$ be the graph given by \cref{con:HG}. 
	By \cref{prop:no-approx-hom}, there is no $\varepsilon$-approximate homomorphism from $G^{\star}$ to any $F$-hom-free graph on at most $2^{cm/n}$ vertices. Note that $2^{cm/n} = 2^{\Omega(n^{1/g})} \geq 2^{(1/\varepsilon)^{1/g}}$ (provided that $\varepsilon$ is small enough). 

	It remains to show that $G^{\star}$ is $H$-hom-free. So suppose for contradiction that there is a homomorphism $\psi:H \to G^\star$. 
	For a vertex $u \in V(G)$, let $V_u \coloneqq \{u\} \times \{1,\dots,\Delta\}^m$ denote the blowup set of $u$ in $G^{\star}$ (recall \cref{con:HG}).  
	Let $X$ be the set of all $u \in V(G)$ such that $V_u \cap \psi(V(H)) \neq \emptyset$; so $\abs X \leq \abs H = h \leq g$. 
    Let $F_1,\dots,F_m$ be an enumeration of the $F$-copies in $G$. We consider the $F$-copies in $G$ which intersect $X$ in at least 2 vertices; without loss of generality, these are 
   $F_1,\dots,F_s$. As we saw above, the hypergraph 
    $\{V(F_i) \cap X : 1 \leq i \leq s\}$ is a hyperforest. 
    Let $T$ be the $F$-forest defined as follows: $T$ has $s$ copies of $F$, denoted $F'_1,\dots,F'_s$; $V(T)$ consists of the set $X$ and the sets $V(F'_i) \setminus X$ (for $i = 1,\dots,s$), which are pairwise-disjoint; and $T[V(F'_i) \cap X] = G[V(F_i) \cap X]$ as labelled graphs.\footnote{Namely, the sets $V(F_i)\setminus X$ may intersect (in $G$), but we make these sets disjoint in $T$ to make sure that $T$ is an $F$-forest. We could avoid this technicality by increasing $g$ to $\binom{h}{2}(f-2)$, which would ensure that $G[\bigcup_{i=1}^s V(F_i)]$ is an $F$-forest.}
    Then $T$ is homomorphic to $G$ via the homomorphism which maps $F'_i$ isomorphically to $F_i$ for every $i \in [s]$. As $G$ is homomorphic to $F$, we get that $T$ is homomorphic to $F$.
    
	Let $T^{\star}$ be the graph obtained by applying \cref{con:HG} to $T$. For a vertex $(v,\bx) \in V(G^{\star})$, let $\phi(v,\bx)$ be the result of projecting $\bx$ on the first $s$ coordinates, i.e., if $\bx = (\bx_1,\dots,\bx_m)$, then $\phi(v,\bx) = (v,(\bx_1,\dots,\bx_s))$. If $\mathbf{v} \in \psi(V(H))$ then $v \in X$, so $\phi(v,\bx) \in V(T^{\star})$. We claim that $(\phi \circ \psi) : V(H) \rightarrow V(T^{\star})$ is a homomorphism. Indeed, let $u_0,v_0$ be adjacent vertices in $H$, and let $(u,\bx)=\psi(u_0)$ and $(v,\by) = \psi(v_0)$. These two vertices are adjacent in $G^\star$ since $\psi$ is a homomorphism. Hence, $uv \in E(G)$ by the definition of $G^{\star}$. Also, $u,v \in X$, so there exists 
    $1 \leq k \leq s$ such that $uv \in E(F_k)$. Let $(u,i),(v,j)$ be the labels of the edge $uv$ in $F_k$. By the definition of $G^{\star}$, the fact that $(u,\bx)$ and $(v,\by)$ are adjacent means that $\bx_k = i$ and $\by_k = j$. But $T^{\star}$ has the same adjacency criterion, meaning that $\phi(u,\bx)$ and $\phi(v,\by)$ are adjacent in $T^{\star}$. Since this holds for arbitrary $u_0 v_0 \in E(H)$, we conclude that $\phi \circ \psi$ is indeed a homomorphism $H \to T^\star$. 
    Thus, $H$ is homomorphic to $T^\star$ for an $F$-forest $T$ with $T \rightarrow F$.
    This contradicts the assumption of the theorem and concludes the proof.
	\end{proof}

\end{document}
